% Commandes spécifiques ou pour la mise en forme

\makeatletter
\newcommand\xleftrightarrow[2][]{%
  \ext@arrow 9999{\longleftrightarrowfill@}{#1}{#2}}
\newcommand\longleftrightarrowfill@{%
  \arrowfill@\leftarrow\relbar\rightarrow}
\makeatother

\hypersetup{
    linkcolor=linkcolor, % couleur des liens internes (acronymes, toc, refs…)
    citecolor=linkcolor, % couleur des citations
    urlcolor=linkcolor   % couleur des URLs
}


%\newacronym{LL}{LL}{Lieb-Liniger}
%\newacronym{NS}{NS}{Schrödinger non linéaire}
%\newacronym{GP}{GP}{Gross–Pitaevskii}
%\newacronym{GGE}{GGE}{Generalized Gibbs Ensemble}
%\newacronym{BA}{BA}{Bethe Ansatz}
%\newacronym{FBA}{FBA}{Fermions Bethe Ansatz}
%\newacronym{TBA}{TBA}{Thermodynamique de Bethe Ansatz}
%\newacronym{qBEC}{qBEC}{quasi-condensat de Bose-Einstein}
%\newacronym{GHD}{GHD}{Hydrodynamique Généralisée}
%\newacronym{LDA}{LDA}{Approximation de Densité Locale}

\newacronym{LL}{LL}{Lieb–Liniger}
\newacronym{NS}{NLS}{Nonlinear Schrödinger}
\newacronym{GP}{GP}{Gross–Pitaevskii}
\newacronym{GGE}{GGE}{Generalized Gibbs Ensemble}
\newacronym{BA}{BA}{Bethe Ansatz}
\newacronym{FBA}{FBA}{Fermionic Bethe Ansatz}
\newacronym{TBA}{TBA}{Thermodynamic Bethe Ansatz}
\newacronym{qBEC}{qBEC}{quasi Bose–Einstein Condensate}
\newacronym{GHD}{GHD}{Generalized Hydrodynamics}
\newacronym{LDA}{LDA}{Local Density Approximation}
\newacronym{DMD}{DMD}{Dispositif de Micromirroirs Digitaux}
\newacronym{PMO}{PMO}{Piège Magnéto-Optique}
\newacronym{DC}{DC}{Direct Current}
\newacronym{AC}{AC}{Alternating Current}
\newacronym{SYRTE}{SYRTE}{Systèmes de Référence Temps-Espace} 
\newacronym{DFB}{DFB}{Distributed Feedback}
\newacronym{DBR}{DBR}{Distributed Bragg Reflector}
\newacronym{TA}{TA}{Tapered Amplifier}
\newacronym{AOM}{AOM}{Acousto-Optic Modulator}
\newacronym{PBS}{PBS}{Polarizing Beam Splitter}
\newacronym{LCF}{LCF}{Laboratoire Charles Fabry}
\newacronym{C2N}{C2N}{Centre de nanosciences et de nanotechnologies}
\newacronym{BCB}{BCB}{benzocyclobuténe}
\newacronym{SiC}{SiC}{silicon carbide}
\newacronym{Au}{Au}{aurum}
%\newacronym{AOM}{AOM}{modulateur acousto-optique}
\newacronym{CCD}{CCD}{Charge-Coupled Device}
\newacronym{TF}{TF}{Thomas-Fermi}
\newacronym{CESQ}{CESQ}{Centre européen de sciences quantiques}


\setcounter{secnumdepth}{3}  % numérotation jusqu’aux subsubsections
\setcounter{tocdepth}{2}     % affichage de la ToC jusqu’aux subsections


